\scanpage{85}
$(\unsure{E|\mathcal E}_{p(x)})_{/X}$ \deleted{$\unsure{\simeq E_x|\simeq\widetilde E_x}$},
\deleted{mais en \unsure{astreignant|assignant}} qui sont
cartésiennes \unsure{quand|qd} $X$ parcourt une fibre
(``petit site'') $E_x$. De façon précise

\[
\widetilde E^\circ \approxeq \underline{\Gamma}^{\mathrm{vert\text{-}cart}}
(\deleted{\unsure{B}}E,\ \deleted{\unsure{E}}X\mathrel{\unsure{\leadsto|\sim}}
{(\overline E)_{p(X)/X}}^\circ),
\]

où l'exposant ``vert-cart'' signifie qu'on
se borne aux sections qui sont cartésiennes
au dessus des flèches verticales. Utilisant
une formule d'associativité évidente, le
deuxième \unsure{ms|membre} s'interprète comme

\[
\underline{\Gamma}\bigl(B,\ x\leadsto
\underbrace{\underline{\Gamma}^{\mathrm{cart}}\bigl(E_x,\ X\leadsto
(\overline E_x)^\circ_{\!/X\unsure{\rbrace|3}}\bigr)}_{\displaystyle\widetilde E_x^\circ=(\overline E)_x^\circ}
\bigr)
\]

Donc

\[
\widetilde E\simeq\underline{\Gamma}(B,\overline E^\triangle)^\circ
\]

où $\overline E^\triangle$ désigne la catégorie fibrée déduite de
$\overline E$ en passant aux \deleted{\unsure{duales}} \inserted{opposées} des fibres de $\overline E$.

L'inclusion

\[
\underleftarrow{\mathrm{Lim}}\bigl(E,X\leadsto(\overline E)_{p(X)/X}\bigr)
\longrightarrow\overset{\widetilde E}{\overset{\Vert}{\underline{\Gamma}^{\mathrm{vert.cart}}}}
(E,\ldots)
\]

n'est alors autre \inserted{\unsure{mod signis|mod isom}} que l'inclusion

\[
\unsure{[|\underline{\Gamma}}\ \underleftarrow{\mathrm{Lim}}\,(B,\overline E^\triangle)^\circ
\simeq\underleftarrow{\mathrm{Lim}}\,(B,\overline E)^{\deleted{\unsure{\sharp|*|\circ|\triangle}}}
\subset\underline{\Gamma}(B,\overline E^\triangle)^\circ
\]

% archive sheet number: 84
